\subsubsection{}

Vi modellerer ionene som punktladninger med ladning $+e$ og $-e$. Coulombs lov gir kraften mellom ladningene:
\begin{align}
\vec{F}_c = k_e \frac{qQ}{r^2} \,\hat{e}_r
\end{align}

Den potensielle energien kan finnes fra:
\begin{align}
E_p(b) - E_p(a) = - \int_a^b \vec{F} \cdot d\vec{s}
\end{align}

Vi velger referansepunkt $a = \infty$, hvor $E_p(\infty) = 0$, og $b = r$:
\begin{align}
E_p(r) = - \int_{\infty}^{r} \vec{F}_c \cdot d\vec{s}
\end{align}

Siden bevegelsen skjer i radiell retning, kan vi sette $d\vec{s} = \hat{e}_r \, dr$:
\begin{align}
E_p(r) = - \int_{\infty}^{r} k_e \frac{qQ}{r^2} \,\hat{e}_r \cdot \hat{e}_r \, dr
\end{align}

Siden $\hat{e}_r \cdot \hat{e}_r = 1$, får vi:
\begin{align}
E_p(r) = - k_e qQ \int_{\infty}^{r} \frac{1}{r^2} \, dr
\end{align}

Integralet gir:
\begin{align}
E_p(r) = - k_e qQ \left[ -\frac{1}{r} \right]_{\infty}^{r} = k_e \frac{qQ}{r}
\end{align}

Setter inn verdier:
\begin{align}
qQ = -(1.602 \cdot 10^{-19}\,\text{C})^2
\end{align}

\begin{align}
r = 2.36 \cdot 10^{-10}\,\text{m}, \quad k_e = 8.99 \cdot 10^9\,\text{Nm}^2/\text{C}^2
\end{align}

\begin{align}
E_p(r) = \frac{(8.99 \cdot 10^9)\cdot (-(1.602 \cdot 10^{-19})^2)}{2.36 \cdot 10^{-10}}
\approx -9.78 \cdot 10^{-19}\,\text{J}
\end{align}

Omregnet til elektronvolt:
\begin{align}
E_p \approx -6.1\,\text{eV}
\end{align}

\newpage